\documentclass[11pt,a4paper]{article}
\usepackage[utf8]{inputenc}
\usepackage[T1]{fontenc}
\usepackage[english]{babel}
\usepackage[margin=1in]{geometry}
\usepackage{amsmath,amssymb}
\usepackage{graphicx}
\usepackage[most]{tcolorbox}
\usepackage{booktabs}
\usepackage{array}
\usepackage{longtable}
\usepackage{enumitem}
\usepackage{fancyhdr}
\usepackage{xcolor}
\definecolor{unalblue}{RGB}{0,56,101}
\definecolor{boxbg}{RGB}{245,248,250}
\newtcolorbox{resultbox}[1][]{colback=boxbg,colframe=unalblue,title=#1,fonttitle=\bfseries}
\pagestyle{fancy}
\fancyhf{}
\rhead{\thepage}
\usepackage[hidelinks]{hyperref}
\title{Local Research Agent Document Report}
\author{Local Linux Research Agent}
\date{\today}
\begin{document}
\maketitle
\thispagestyle{fancy}

\section\emph{{Chat History & LLM Inferences}
\begin{itemize}
\noindent \textbf{1.} user: FM\\\n\noindent \textbf{2.} agent: I found 1 local document related to "FM".
This first prototype uses local document metadata and filenames. Enable Ollama plus text extraction/embeddings for deeper inference over PDF contents.

\item 03SpectrumSensingFM0 (03SpectrumSensingFM0.pdf, 488294 bytes, sha256 3e283c6d13cf610d4276a7e64711ba4441938807abb99de3f772e2cc049cb1b2, type pdf)\\\n\noindent \textbf{3.} user: [Deep] Cuáles son las 6 etapas del pipeline de monitoreo FM y qué hace cada una?"\\\n\noindent \textbf{4.} agent: [Deep Analysis of 03SpectrumSensingFM0]

Heuristic Analysis: The document mentions "Cuáles son las 6 etapas del pipeline de monitoreo FM y qué hace cada una?"" in several sections. Key findings related to this query include: Uncertainty Budget and Reporting
All measurements intended for regulatory compliance use shall be accompanied by a documented
uncertainty  evaluation  in  accordance  with  the  Guide  to  the  Expression  of  Uncertainty  in  Mea-
surement (GUM) [?].  The uncertainty budget shall account for all known sources of random and
systematic error, shall be traceable to the calibration hierarchy defined above, and shall be reported
alongside the measurand value in all compliance records and evidence archives.
Uncertainty Classification, the associated Uncertainty-Aware Compliance Decisions, and Peri-
odic Re-Evaluation Requirements are detailed in Appendix
\item Node-Level DSP Pipeline for FM Compliance Monitoring
The compliance monitoring pipeline transforms raw complex baseband observations into reported
measurands, decision outcomes, and supporting evidence.  However, the technical rigor of such a
pipeline implementation is often challenged by the hardware-software interface.
The principal processing stages are shown in Fig. 2 and described below.
Wideband
Acquisition
Conditioning
Channeli-
zation \&
Tracking
Per-Channel
Demodulation
Identification
Compliance
Measurand
Estimation
Rule
Evaluation
\& Alerting
Reporting
\& Evidence
Archiving
Input:
Raw I/Q
Samples
Output:
Compliance Reports
\& Evidence Logs
Figure 2:  Implemented SDR-based FM broadcast compliance monitoring pipeline.  Raw I/Q sam-
ples are processed through six stages, from wideband acquisition and conditioning to reporting and
evidence archiving, to enable automated regulatory assessment.
– Wideband Acquisition \& Conditioning:  Raw I/Q samples are acquired from the SDR front end
together with receiver-configuration and timing metadata.  This stage establishes the wide-
band observation used by all downstream analyses.
Its principal functions are:
ˆ  Hardware  configuration:   initialization  of  SDR  operating  parameters,  including  center
frequency, sample rate, gain state, and front-end amplification settings;
ˆ  Sample  streaming:  transfer of raw complex baseband samples from the SDR device to
the embedded host processor through the USB interface;
ˆ  Buffer  management:  temporary storage and flow control to reduce the risk of sample
loss during sustained acquisition;
ˆ  Metadata association:  tagging of captured data with timing, configuration, and calibration-
state information;
ˆ  Signal  conditioning:   preliminary  preparation  of  the  wideband  observation,  including
DC-offset mitigation, optional digital frequency translation, and other conditioning steps
required before channelization and spectral analysis.
– Channelization, Isolation \& Tracking:  The conditioned wideband observation is analyzed to de-
tect active FM broadcast channels within the monitored span.  Detected channels are isolated
and tracked over time to support stable downstream measurement.
Its principal functions are:
8; at 75 kHz peak deviation of a 1 kHz sine tone, with the 19 kHz stereo pilot present at
its nominal level.  This definition shall be recorded in the measurement procedure.
ˆ  Sub-carrier  verification.   The  19  kHz  pilot  tone  and  57  kHz  RDS  sub-carrier  shall  be
independently  verified  in  magnitude.   Results  shall  be  reported  relative  to  the  main
carrier level and stored as supporting evidence.
– Channel Occupancy.ˆ  Detection threshold. The occupancy threshold shall be a site-calibrated
offset above the locally verified noise floor, derived from the Tier 2 noise-floor verifica-
tion required by Table 3 precondition (2).  A fixed default offset shall not be used for
Conditional Compliance-Grade reporting.
ˆ  Statistical  reporting.   Occupancy shall  be  reported  as  the  fraction of  observation  win-
dows in which the channel is declared occupied under the documented detection rule,
together with a confidence interval derived from the characterised false-alarm and missed-
detection  probabilities.   The  reporting  model  shall  be  consistent  with  the  statistical
probability framework required by Table 3 precondition (3).
ˆ  Hysteresis  and  dwell.  Temporal hysteresis shall be applied to prevent false occupancy
triggers from intermittent burst signals.  The hysteresis parameters shall be fixed in the
version-controlled processing configuration per Table 3 precondition (4).
Estimation Pipeline of Measurands
Reference Carrier-Frequency Estimation Procedure.    The processing pipeline is specified
for estimating the carrier centre frequency of an FM broadcast signal from complex baseband sam-
ples acquired with a HackRF One software-defined radio receiver.  The implementation is organised
as six sequential stages.  The pipeline is explicitly designed to retain intermediate quantities so that
spectral estimates, candidate selection, carrier estimates, and confidence scores remain auditable
and can be reinterpreted if calibration metadata are updated.
– Stage 1:  Acquisition Configuration.  The acquisition stage constructs a validated hardware con-
figuration for one HackRF capture session.  In the notebook implementation, the receiver is
configured with
f
LO
= f
target
+ f
off
,   f
off
=
f
s
4
,
so the observed FM channel is offset from the DC bin by one quarter of the sample rate.  The
configuration validator additionally requires
|f
off
|≥ 300 kHz,  |f
off
| + 150 kHz <
f
s
2
,
so that the target FM allocation remains well separated from the DC region and inside the
usable complex baseband span.
The user-facing HackRF receive controls are:
ˆ  RF amplifier state ampenable (internal +14 dB front-end amplifier),
ˆ  IF gain ifgaindb in 8 dB steps over 0 to 40 dB,
ˆ  baseband gain basebandgaindb in 2 dB steps over 0 to 62 dB.
20; Reference Requirements for Full Compliance-Grade Deployments
The following requirements apply to each measurand when compliance-grade or Conditional Compliance-
Grade reporting is sought.  Requirements for measurands classified as Unsupported on the baseline
HackRF One platform (field strength, ACLR) are scoped to upgraded or augmented platforms as
defined in the Upgrade Paths section; they do not imply that these measurands are achievable on
the unaugmented HackRF One.  All requirements are in addition to, not in substitution for, the
preconditions listed in Table 3.
– Carrier Frequency Error and Stability.ˆ  Reference clock traceability.  The SDR master clock
shall be disciplined by a GPS-Disciplined Oscillator (GPSDO) or equivalent traceable
reference providing a fractional frequency accuracy of ±1× 10
−9
or better.  An external
Rubidium  oscillator  is  acceptable  only  when  it  is  itself  GPS-steered  or  carries  a  cur-
rent  calibration  certificate  traceable  to  a  national  frequency  standard;  a  free-running
Rubidium standard does not satisfy this requirement.
ˆ  Estimation  algorithm.   Carrier  centre  frequency  shall  be  estimated  using  the  Stage  5
pipeline  specified  in  the  Estimation  Pipeline  section  (parabolic  interpolation  primary,
power centroid and edge midpoint as supporting estimators).  The power-weighted cen-
troid  may  be  used  as  the  primary  estimator  only  when  the  peak-parabolic  method  is
flagged invalid by the Stage 5 validity logic.
ˆ  Guard-band decision logic.  The Guard-Band rule (Eq. 4) shall be applied using the 2 kHz
carrier-frequency  tolerance  specified  by  FCC  Part  73  or  the  applicable  jurisdictional
limit.   If |f
cal
− f
target
| + U ≤ T ,  the  result  is  compliant;  if |f
cal
− f
target
|− U > T ,
non-compliant; otherwise uncertain and subject to the escalation procedure in Table 4.
ˆ  Uncertainty-driven  escalation.   As  specified  in  Table  4,  results  in  the  uncertain  zone
shall trigger a repeated measurement cycle with 4× extended averaging.  If the result
remains uncertain after three attempts, it shall be logged as UNRESOLVED and escalated
to operator review.
ˆ  Reference drift monitoring.  The GPSDO lock state and 1 PPS timing validity shall be
recorded continuously.  Any measurement acquired during a GPSDO unlock event shall
be flagged UNCALIBRATED
REFERENCE and classified as screening-grade.  Temperature of
the SDR front-end shall be logged as a diagnostic field; if the node operates without a
disciplined reference (e.g., during holdover), the estimated frequency drift based on the
TCXO drift profile shall be recorded as an additional Type B uncertainty contribution.
– Calibrated Received Power.ˆ  Reference  plane  traceability.   The  receive  chain  shall  undergo
Tier 1 laboratory calibration establishing absolute gain accuracy at the defined reference
plane (antenna port or SDR input bulkhead).  Tier 2 field verification shall confirm that
the node has not drifted beyond acceptable limits relative to the Tier 1 state;  Tier 2
alone does not establish absolute traceability.
ˆ  Shared-risk  decision  logic.  Expanded uncertainty U  (k = 2) shall be computed in real
time.  Values falling within the U -interval of the applicable ERP or received-power limit
shall  be  flagged  for  verification  using  a  calibrated  mobile  unit  or  a  compliance-grade
co-located node, per Table 4.
ˆ  Integration method.  Calibrated received power shall be computed as the RMS-weighted
PSD integral over the defined analysis bandwidth B (nominally 200 kHz for FM broad-
cast), after gain-state and frequency-response correction.
18.\\
\end{itemize}

\section}{Indexed Local Documents}
This report was generated from the local document mount configured by \texttt{DOCS\_DIR}.

\section*{1. 03SpectrumSensingFM0}
\textbf{File:} 03SpectrumSensingFM0.pdf\\
\textbf{SHA-256:} \texttt{3e283c6d13cf610d4276a7e64711ba4441938807abb99de3f772e2cc049cb1b2}\\
\textbf{Size:} 488294 bytes\\
\textbf{Indexed:} 2026-08-14T14:14:10.171Z

\end{document}
